% 第 2 章 Licensing and Resource Requirements
\bichapter{许可与资源需求}{Licensing and Resource Requirements}
\label{chap:licensing}

本章介绍 IC Validator 工具的设置说明。请确认目录与文件齐全、运行环境满足 IC Validator 要求，并已正确配置许可信息。

许可与资源需求包括以下各节：
\begin{itemize}
  \item 许可方案（Licensing Schemes）
  \item 许可等待（License Waiting）
  \item 需要哪些 IC Validator 许可证（Which IC Validator Licenses Do I Need?）
  \item 资源需求（Resource Requirements）
\end{itemize}

关于 Synopsys 许可及 IC Validator 安装的更多信息，请参阅：
\begin{itemize}
  \item \textit{Synopsys Licensing QuickStart Guide}，见\\
        \url{http://www.synopsys.com/Support/LI/Licensing/Pages}
  \item IC Validator 安装信息，见\\
        \url{http://www.synopsys.com/Support/LI/Installation/Pages/default.aspx}
\end{itemize}

% ============================================================
\section{许可方案}
\subsection*{Licensing Schemes}

IC Validator 工具有三种许可方案，运行时使用其中之一：
\begin{itemize}
  \item Apex 许可（Apex Licensing）
  \item Elite 许可（Elite Licensing）
  \item Base 许可（Base Licensing）
\end{itemize}

Base 许可为默认方案。若作业未使用许可等待，也未指定 \cmd{-lic\_apex}、\cmd{-lic\_elite} 或 \cmd{-lic\_base} 命令行选项，则 IC Validator 首先检查该次运行所需的最低档（tier）许可证；若不可用，再检查更高档。若运行可使用任一档，则先检查 Base；Base 不可用时检查 Elite，再检查 Apex。若均不可用，作业以许可拒绝（license denied）错误结束。

三种许可方案可同时作为系统上的可用方案存在，但一次 IC Validator 作业只使用其中一种，不能混用。

若多种许可证均可用，可用 \cmd{-lic\_apex}、\cmd{-lic\_elite} 或 \cmd{-lic\_base} 为特定作业选择许可方案。更多信息见第~\ref{chap:basics}~章「IC Validator 命令行选项」。

% ============================================================
\section{许可等待}
\subsection*{License Waiting}

IC Validator 按 Synopsys Common Licensing（SCL）标准指南支持许可等待（license waiting）。

在启用许可等待时运行 IC Validator：
\begin{itemize}
  \item 若仅安装了 Apex 许可，作业排队等待 Apex 许可证。
  \item 若仅安装了 Elite 许可，作业排队等待 Elite 许可证。
  \item 若仅安装了 Base 许可，作业排队等待 Base 许可证。
  \item 若安装了多种许可方案，作业按 Base、Elite、Apex 的顺序，排队等待其中第一种已安装的方案；其余方案被忽略。
  \begin{itemize}
    \item 若使用 \cmd{-lic\_base}，作业排队等待 Base，忽略 Apex 与 Elite。
    \item 若使用 \cmd{-lic\_elite}，作业排队等待 Elite，忽略 Apex 与 Base。
    \item 若使用 \cmd{-lic\_apex}，作业排队等待 Apex，忽略 Elite 与 Base。
  \end{itemize}
\end{itemize}

\subsection{用户控制}
\subsubsection*{User Controls}

当提交作业时并非全部多核（multicore）许可证都可用，IC Validator 提供如下用户控制：

\begin{description}[style=nextline, leftmargin=1.2em]
  \item[\cmd{License\_wait\_duration}]
    指定工具等待许可证变为可用的最长时间。
  \item[\cmd{License\_mode}]
    决定启动时并非全部多核许可证都可用时的行为。
\end{description}

\begin{itemize}
  \item \cmd{run\_available}：作业以当前可用的部分许可证运行。
  \begin{itemize}
    \item 若同时设置了 \cmd{-license\_wait\_duration} 且当时没有任何许可证，作业排队直至有非零数量的许可证可用，然后用这些许可证运行。
  \end{itemize}
  \item \cmd{run\_and\_add}：作业以当前可用的部分许可证运行，随后继续排队并在许可证变为可用时追加签出，尽量达到所请求的全部许可证数量。
  \begin{itemize}
    \item 若同时设置了 \cmd{-license\_wait\_duration} 且当时没有任何许可证，作业排队直至有非零数量可用并以此启动；之后继续排队追加，尽量达到全部请求数量。
  \end{itemize}
  \item \cmd{exit}：若完整数量不可用，作业退出。
  \begin{itemize}
    \item 若同时设置了 \cmd{-license\_wait\_duration}，作业排队并逐步获取许可证；若在时限内达到完整数量则启动，否则退出。
  \end{itemize}
  \item \cmd{exit\_and\_retry}：若完整数量不可用，作业退出、释放已获取的许可证并重试。
  \begin{itemize}
    \item 若同时设置了 \cmd{-license\_wait\_duration}，作业排队并逐步获取；若在时限内达到完整数量则启动；否则退出、释放已获取许可证并重试。
  \end{itemize}
\end{itemize}

% ============================================================
\section{许可缓存}
\subsection*{License Caching}

IC Validator 在获知各许可服务器上安装了哪些许可密钥后会缓存这些信息。在后续运行中，缓存信息可能用于调整作业使用的许可服务器列表顺序，或调整检查各许可方案的顺序，从而可能缩短成功签出许可证的时间。

缓存信息会周期性过期，并随工具进一步了解各服务器上的可用许可证而更新。当工具改变许可服务器或许可方案的优先级时，会记录到 \cmd{run\_details/licmsg}。

默认情况下，许可缓存在下列目录：
\begin{lstlisting}
$XDG_CACHE_HOME/Synopsys/icv/license
\end{lstlisting}

若未设置 \cmd{\$XDG\_CACHE\_HOME}，则许可缓存在：
\begin{lstlisting}
~/.cache/Synopsys/icv/license
\end{lstlisting}

可设置环境变量 \cmd{ICV\_LICENSE\_CACHE\_DIR} 覆盖许可缓存位置。

使用命令行选项 \cmd{-nlc} 可禁用许可缓存。

% ============================================================
\section{需要哪些 IC Validator 许可证}
\subsection*{Which IC Validator Licenses Do I Need?}

runset 中使用的 IC Validator 功能决定所需许可证。以下各节说明所需许可：
\begin{itemize}
  \item Apex 许可（Apex Licensing）
  \item Elite 许可（Elite Licensing）
  \item Base 许可（Base Licensing）
  \item 按功能的许可需求（Licensing Requirements by Function）
  \item 按技术的许可需求（Licensing Requirements by Technology）
  \item 多核许可（Multicore Licensing）
  \item 令牌许可（Token Licensing）
  \item 超线程（Hyperthreading）
\end{itemize}

\subsection{Apex 许可}
\subsubsection*{Apex Licensing}

Apex 许可使用下列密钥：
\begin{itemize}
  \item \cmd{ICValidator-Manager-Apex}
  \item \cmd{ICValidator2-GeometryEngine}
  \item \cmd{ICValidator2-CompareEngine}
\end{itemize}

每次运行 IC Validator 都会签出 \cmd{ICValidator-Manager-Apex}；该密钥采用分档许可：高档 Apex、中档 Elite、低档 Base，具体取决于所用技术或功能。

\subsection{Elite 许可}
\subsubsection*{Elite Licensing}

Elite 许可使用下列密钥：
\begin{itemize}
  \item \cmd{ICValidator-Manager}
  \item \cmd{ICValidator2-GeometryEngine}
  \item \cmd{ICValidator2-CompareEngine}
\end{itemize}

每次运行 IC Validator 都会签出 \cmd{ICValidator-Manager}；同样按技术或功能采用 Apex / Elite / Base 分档。

\subsection{Base 许可}
\subsubsection*{Base Licensing}

Base 许可使用下列密钥：
\begin{itemize}
  \item \cmd{ICValidator-Manager-2020}
  \item \cmd{ICValidator2-GeometryEngine}
  \item \cmd{ICValidator2-CompareEngine}
\end{itemize}

每次运行 IC Validator 都会签出 \cmd{ICValidator-Manager-2020}；同样按技术或功能采用 Apex / Elite / Base 分档。

\subsection{按技术的许可需求}
\subsubsection*{Licensing Requirements by Technology}

技术分组可能需要一个或多个许可证，如表~\ref{tab:icv-tech-license} 所示。

\begin{longtable}{@{}p{0.48\textwidth} p{0.44\textwidth}@{}}
\caption{技术许可需求（Technology Requirements）}\label{tab:icv-tech-license}\\
\toprule
\textbf{技术} & \textbf{所需许可} \\
\midrule
\endfirsthead
\toprule
\textbf{技术} & \textbf{所需许可} \\
\midrule
\endhead
\bottomrule
\endfoot
DRC 功能 & 各档均可用 \\
LVS 提取功能 & 各档均可用 \\
LVS 比较功能 & 各档均可用。使用 \cmd{ICValidator2-CompareEngine} \\
Metal fill 功能 & 各档均可用 \\
Antenna & 各档均可用 \\
与 IC Compiler II / Fusion Compiler 的 Fusion（InDesign） & 各档均可用 \\
Netlist Visualizer & 各档均可用 \\
Layer Debugger & 各档均可用 \\
LVS Short Finder & 各档均可用 \\
DRC Heat Map and Balanced Error Distribution & 各档均可用 \\
IC Validator Launch & 各档均可用 \\
错误数据库（PYDB）查询 & 各档均可用 \\
PERC & 各档均可用 \\
Live DRC - Custom & 各档均可用 \\
Live DRC - Digital & 各档均可用 \\
16-Threaded Commands & 各档均可用 \\
VUE & 各档均可用。使用 \cmd{ICValidator2-GeometryEngine}\textsuperscript{1} \\
NetTran 及其他实用程序 & 各档均可用。使用 \cmd{ICValidator2-GeometryEngine}\textsuperscript{2} \\
Pattern matching 功能 & Elite（\cmd{ICValidator-Manager}）或 Apex（\cmd{ICValidator-Manager-Apex}） \\
Multiple-Patterning Technology（MPT）功能 & Elite（\cmd{ICValidator-Manager}）或 Apex（\cmd{ICValidator-Manager-Apex}） \\
Explorer DRC & Elite（\cmd{ICValidator-Manager}）或 Apex（\cmd{ICValidator-Manager-Apex}） \\
Explorer LVS & Elite（\cmd{ICValidator-Manager}）或 Apex（\cmd{ICValidator-Manager-Apex}） \\
Elastic CPU Management & Elite（\cmd{ICValidator-Manager}）或 Apex（\cmd{ICValidator-Manager-Apex}） \\
Blackbox DRC & Elite（\cmd{ICValidator-Manager}）或 Apex（\cmd{ICValidator-Manager-Apex}） \\
3DIC & Elite（\cmd{ICValidator-Manager}）或 Apex（\cmd{ICValidator-Manager-Apex}） \\
Unified fill & 部分子组件限于 Elite \textbar\ Apex；另有部分子组件限于 Apex \\
ML resource estimation & Apex（\cmd{ICValidator-Manager-Apex}） \\
\end{longtable}

\noindent\small
\textsuperscript{1}仅检查是否存在有效许可证，不签出许可证。\\
\textsuperscript{2}仅检查是否存在有效许可证，不签出许可证。

\subsection{按功能的许可需求}
\subsubsection*{Licensing Requirements by Function}

个别函数可能需要一个或多个许可证，如表~\ref{tab:icv-func-license} 所示。

\begin{longtable}{@{}>{\raggedright\arraybackslash\ttfamily}p{0.46\textwidth} p{0.46\textwidth}@{}}
\caption{功能许可需求（Functions Requirements）}\label{tab:icv-func-license}\\
\toprule
\textbf{\textrm{函数（选项）}} & \textbf{\textrm{所需许可}} \\
\midrule
\endfirsthead
\toprule
\textbf{\textrm{函数（选项）}} & \textbf{\textrm{所需许可}} \\
\midrule
\endhead
\bottomrule
\endfoot
adjust\_color\_graph() &
  \textrm{Elite（\cmd{ICValidator-Manager}）或 Apex（\cmd{ICValidator-Manager-Apex}）} \\
external2( &
  \textrm{各档均可用} \\
\quad look\_thru &
  \textrm{Apex（\cmd{ICValidator-Manager-Apex}）} \\
\quad blocking\_layer &
  \textrm{Apex（\cmd{ICValidator-Manager-Apex}）} \\
) & \\
fill\_to\_text() &
  \textrm{Elite（\cmd{ICValidator-Manager}）或 Apex（\cmd{ICValidator-Manager-Apex}）} \\
four\_color() &
  \textrm{Elite（\cmd{ICValidator-Manager}）或 Apex（\cmd{ICValidator-Manager-Apex}）} \\
layout\_text\_to\_net\_property() &
  \textrm{Elite（\cmd{ICValidator-Manager}）或 Apex（\cmd{ICValidator-Manager-Apex}）} \\
net\_color\_check() &
  \textrm{Elite（\cmd{ICValidator-Manager}）或 Apex（\cmd{ICValidator-Manager-Apex}）} \\
net\_select\_error() &
  \textrm{Elite（\cmd{ICValidator-Manager}）或 Apex（\cmd{ICValidator-Manager-Apex}）} \\
not\_within\_distance\_by\_property() &
  \textrm{Elite（\cmd{ICValidator-Manager}）或 Apex（\cmd{ICValidator-Manager-Apex}）} \\
optional\_pattern\_markers() &
  \textrm{Elite（\cmd{ICValidator-Manager}）或 Apex（\cmd{ICValidator-Manager-Apex}）} \\
pattern\_extract() &
  \textrm{Elite（\cmd{ICValidator-Manager}）或 Apex（\cmd{ICValidator-Manager-Apex}）} \\
pattern\_learn() &
  \textrm{Elite（\cmd{ICValidator-Manager}）或 Apex（\cmd{ICValidator-Manager-Apex}）} \\
pattern\_library\_compare() &
  \textrm{Elite（\cmd{ICValidator-Manager}）或 Apex（\cmd{ICValidator-Manager-Apex}）} \\
pattern\_library\_lock() &
  \textrm{Elite（\cmd{ICValidator-Manager}）或 Apex（\cmd{ICValidator-Manager-Apex}）} \\
pattern\_library\_merge() &
  \textrm{Elite（\cmd{ICValidator-Manager}）或 Apex（\cmd{ICValidator-Manager-Apex}）} \\
pattern\_library\_read() &
  \textrm{Elite（\cmd{ICValidator-Manager}）或 Apex（\cmd{ICValidator-Manager-Apex}）} \\
pattern\_match() &
  \textrm{Elite（\cmd{ICValidator-Manager}）或 Apex（\cmd{ICValidator-Manager-Apex}）} \\
pattern\_predict() &
  \textrm{Elite（\cmd{ICValidator-Manager}）或 Apex（\cmd{ICValidator-Manager-Apex}）} \\
pattern\_predict\_marker() &
  \textrm{Elite（\cmd{ICValidator-Manager}）或 Apex（\cmd{ICValidator-Manager-Apex}）} \\
pex\_color\_layer\_map() &
  \textrm{Elite（\cmd{ICValidator-Manager}）或 Apex（\cmd{ICValidator-Manager-Apex}）} \\
pex\_qtf\_layer\_map() &
  \textrm{Elite（\cmd{ICValidator-Manager}）或 Apex（\cmd{ICValidator-Manager-Apex}）} \\
pex\_qtf\_layers() &
  \textrm{Elite（\cmd{ICValidator-Manager}）或 Apex（\cmd{ICValidator-Manager-Apex}）} \\
reduce\_four\_color\_graph() &
  \textrm{Elite（\cmd{ICValidator-Manager}）或 Apex（\cmd{ICValidator-Manager-Apex}）} \\
reduce\_three\_color\_graph() &
  \textrm{Elite（\cmd{ICValidator-Manager}）或 Apex（\cmd{ICValidator-Manager-Apex}）} \\
remove\_vias() &
  \textrm{Elite（\cmd{ICValidator-Manager}）或 Apex（\cmd{ICValidator-Manager-Apex}）} \\
three\_color() &
  \textrm{Elite（\cmd{ICValidator-Manager}）或 Apex（\cmd{ICValidator-Manager-Apex}）} \\
two\_color() &
  \textrm{Elite（\cmd{ICValidator-Manager}）或 Apex（\cmd{ICValidator-Manager-Apex}）} \\
unified\_fill( &
  \textrm{各档均可用} \\
\quad different\_size\_fill\_spacing &
  \textrm{Apex（\cmd{ICValidator-Manager-Apex}）} \\
\quad forbidden\_layer &
  \textrm{Apex（\cmd{ICValidator-Manager-Apex}）} \\
\quad merging\_limit\_exclude\_layers &
  \textrm{Apex（\cmd{ICValidator-Manager-Apex}）} \\
\quad merging\_limit\_exclude\_layers\_secondary &
  \textrm{Apex（\cmd{ICValidator-Manager-Apex}）} \\
\quad min\_fill\_region\_concave\_distance &
  \textrm{Apex（\cmd{ICValidator-Manager-Apex}）} \\
\quad projecting\_layer\_shrink &
  \textrm{Apex（\cmd{ICValidator-Manager-Apex}）} \\
\quad remove\_point\_touch &
  \textrm{Apex（\cmd{ICValidator-Manager-Apex}）} \\
) & \\
unified\_fill\_region( &
  \textrm{各档均可用} \\
\quad deangle\_max\_deviation &
  \textrm{Apex（\cmd{ICValidator-Manager-Apex}）} \\
\quad region\_type = UF\_BLOCKAGE\_REGION &
  \textrm{Apex（\cmd{ICValidator-Manager-Apex}）} \\
) & \\
xref\_to\_double\_property() &
  \textrm{Elite（\cmd{ICValidator-Manager}）或 Apex（\cmd{ICValidator-Manager-Apex}）} \\
xref\_to\_double\_property() &
  \textrm{Elite（\cmd{ICValidator-Manager}）或 Apex（\cmd{ICValidator-Manager-Apex}）} \\
xref\_to\_property() &
  \textrm{Elite（\cmd{ICValidator-Manager}）或 Apex（\cmd{ICValidator-Manager-Apex}）} \\
\end{longtable}

\subsection{多核许可}
\subsubsection*{Multicore Licensing}

IC Validator 的 Apex、Elite 与 Base 多核许可需求如下。

\subsubsection{Apex 多核许可}
\paragraph*{Apex Multicore Licensing}

工具可使用一个或多个 \cmd{ICValidator-Manager} 许可证运行（Apex 档对应 \cmd{ICValidator-Manager-Apex}）。每个许可证最多允许使用四个 CPU。例如：
\begin{quote}
100 CPU 作业需要 25 个 Apex（\cmd{ICValidator-Manager-Apex}）许可证
\end{quote}

\subsubsection{Elite 多核许可}
\paragraph*{Elite Multicore Licensing}

工具可使用一个或多个 \cmd{ICValidator-Manager} 许可证运行。每个许可证最多允许使用四个 CPU。例如：
\begin{quote}
100 CPU 作业需要 25 个 Elite（\cmd{ICValidator-Manager}）许可证
\end{quote}

\subsubsection{Base 多核许可}
\paragraph*{Base Multicore Licensing}

工具可使用一个或多个 \cmd{ICValidator-Manager-2020} 许可证运行。每个许可证最多允许使用两个 CPU。例如：
\begin{quote}
100 CPU 作业需要 50 个 Base（\cmd{ICValidator-Manager-2020}）许可证
\end{quote}

\subsection{令牌许可}
\subsubsection*{Token Licensing}

基于令牌（token）的许可描述签出多个许可证的方式。令牌许可与多核许可类似：根据 CPU 数量、按既定比例消耗多个 ICV Manager 许可证。

下列组成部分共同决定一次作业所需的许可证总数：
\begin{quote}
令牌许可证总数 $= ($Per-CPU Flow $\times$ Multicore$) +$ Per-Session Functionality
\end{quote}
\begin{itemize}
  \item \textbf{Per-CPU Flow Functionality}：视作业类型而定，需要一定数量的许可证，并与多核许可相乘。
  \item \textbf{Multicore}：多核结果是 Per-CPU 功能的既定乘数。
  \item \textbf{Per-Session Functionality}：部分功能不与多核相乘；无论多核如何，它们额外消耗固定数量的许可证。
\end{itemize}

\subsection{超线程}
\subsubsection*{Hyperthreading}

IC Validator 签出许可证，直到 CPU 数量等于物理核数。例如：
\begin{quote}
100 CPU 作业需要 25 个 Apex（\cmd{ICValidator-Manager-Apex}）许可证
\end{quote}

\paragraph{超线程示例}
\paragraph*{Hyperthreading Example}

一台启用超线程的 20 核主机呈现为 40 个 CPU，其中 4 个 CPU $= 1$ 个许可证：
\begin{lstlisting}
8 CPUs  (-host_init 8)  : Consumes 2 licenses
12 CPUs (-host_init 12) : Consumes 3 licenses
16 CPUs (-host_init 16) : Consumes 4 licenses
20 CPUs (-host_init 20) : Consumes 5 licenses
24 CPUs (-host_init 24) : Consumes 5 licenses (4 CPUs reserved for Threading)
28 CPUs (-host_init 28) : Consumes 5 licenses (8 CPUs reserved for Threading)
32 CPUs (-host_init 32) : Consumes 5 licenses (12 CPUs reserved for Threading)
36 CPUs (-host_init 36) : Consumes 5 licenses (16 CPUs reserved for Threading)
40 CPUs (-host_init 40) : Consumes 5 licenses (20 CPUs reserved for Threading)
\end{lstlisting}

% ============================================================
\section{资源需求}
\subsection*{Resource Requirements}

初始数据库文件以及 IC Validator 创建与使用的文件决定了系统需求。在对将要检查的设计有一定了解时，可估算物理内存、磁盘空间以及可能的交换空间（swap space）需求。

IC Validator 发行包中的 \cmd{contrib} 目录提供用于分析 IC Validator 性能的脚本。该目录位于发行安装目录顶层。这些脚本不受支持（unsupported）；更多信息见 \cmd{contrib} 目录中的 README 文件。

\subsection{空间需求}
\subsubsection*{Space Requirements}

IC Validator 在由设计数据库生成并存放于 group 目录的 group 文件上执行版图检查。工具首先创建主 group 文件（primary group files）：runset 中每个 \cmd{assign} 函数对应的层各一个文件。每个主 group 文件是描述该层在整个设计中存在情况的数据库，并保持全部层次。随着运行推进，某些检查会生成输出，这些输出也以 group 文件形式写入。例如，group 目录中的文件可能包括：
\begin{lstlisting}
HRCHY.ERROR    inst_indx.dat    pgate.group    toxcont.group
cont.group     met1.group       poly.group     via.group
gate.group     met2.group       psel.group     well.group
inst_data.dat  ngate.group      tox.group
\end{lstlisting}

一般规则：若初始数据库为 GDSII 格式且设计层次风格良好，应为 group 文件创建预留至少相当于初始数据库两到三倍的空间。若设计接近扁平（flat），或存在不希望的数据交互，group 文件创建可能需要初始数据库四到六倍的空间。因此，对 10~MB 的 GDSII 数据库，若设计具有层次性，group 文件至少需要 20--30~MB；扁平设计则至少需要 40--60~MB。

除 group 文件外，IC Validator 运行的其他重要输出还包括运行摘要文件与输出数据库文件。摘要文件为 \cmd{cell.LAYOUT\_ERRORS} 与 \cmd{cell.sum}，均为 ASCII 日志，记录 runset 中所命名 cell 的运行结果。

\cmd{cell.LAYOUT\_ERRORS} 与输出数据库的大小取决于运行中遇到的错误数量。与 group 文件类似，若设计具有层次性，图形输出数据库会小于更扁平的设计：层次设计只需对重复模块内部的错误报告一次，而扁平设计则对模块的每个实例都报告错误。

\subsection{内存使用}
\subsubsection*{Memory Usage}

可通过合理的错误定位策略管理错误输出报告所需的内存：先对设计中的小模块运行 IC Validator 并处理错误；修复后再对更高层次模块运行。这样可使 \cmd{cell.LAYOUT\_ERRORS} 与输出数据库保持合理大小，并自下而上控制错误输出存储所需内存。

IC Validator 完成检查还需要物理内存与交换空间。工具尽量将当前检查所需的全部数据载入内存；物理内存溢出部分使用交换空间。一般来说，一旦需要访问交换空间，性能就会下降，尤其当交换空间分布在流量较大的网络上时。若物理内存足以载入全部数据并在不访问交换空间的情况下完成检查，可获得最大吞吐。

对层次设计，内存约为 GDSII 文件大小三到四倍的机器通常表现良好；扁平设计需要更多硬件。应设置足够的交换空间，使物理内存溢出时不至于耗尽交换空间而导致运行终止。

输出 \cmd{cell.sum} 文件列出完成各项检查所用的内存以及缺页（page faults）次数。若某些检查所需内存超过机器分配量并导致大量数据换出到磁盘，后续运行可能需要更大内存的机器。例~\ref{ex:cell-sum-mem} 给出了内存使用报告示例。下一节「运行时间与内存使用报告」说明摘要文件中报告的运行时间与内存用法。

\begin{lstlisting}[caption={cell.sum 文件中的部分内存使用信息示例},label={ex:cell-sum-mem}]
length_edge() at sf.rs:162
gate_gt_3_5 = length_edge(gate_e, distance > 3.5)
  Function: length_edge
  Inputs: gate_e.edgelayer.0001
  Outputs: gate_gt_3_5.edgelayer.0001
8 unique edge sets written.
  Check time = 0:00:00 User=0.00 Sys=0.00 Mem=0.015 GB
external2() at sf.rs:166
external2(gate_le_1_5, POLY, distance < 0.15, extension = RADIAL,
look_thru = COINCIDENT)
  Comment: "R.POLY.G1: Min Space POLY to Gate for Gate Width <=1.5 must be > 0.15"
  Function: external2
  Inputs: gate_le_1_5.edgelayer.0001, POLY.polygonlayer.0002
WARNING - 6 spacing violations found.
   Check time = 0:00:00 User=0.00 Sys=0.00 Mem=0.016 GB
external2() at sf.rs:169
external2(gate_gt_1_5_le_2_5, POLY, distance < 0.25, extension = RADIAL,
look_thru = COINCIDENT)
  Comment: "R.POLY.G2: Min Space POLY to Gate for 1.5 < Gate Width <=2.5 must be > 0.25"
  Function: external2
  Inputs: gate_gt_1_5_le_2_5.edgelayer.0001, POLY.polygonlayer.0002
WARNING - 6 spacing violations found.
    Check time = 0:00:00 User=0.00 Sys=0.00 Mem=0.016 GB
external2() at sf.rs:172
external2(gate_gt_2_5_le_3_5, POLY, distance < 0.35, extension = RADIAL,
look_thru = COINCIDENT)
  Comment: "R.POLY.G3: Min Space POLY to Gate for 2.5 < Gate Width <=3.5 must be > 0.35"
  Function: external2
  Inputs: gate_gt_2_5_le_3_5.edgelayer.0001, POLY.polygonlayer.0002
WARNING - 3 spacing violations found.
    Check time = 0:00:00 User=0.00 Sys=0.00 Mem=0.016 GB
external2() at sf.rs:175
external2(gate_gt_3_5, POLY, distance < 0.45, extension = RADIAL,
look_thru = COINCIDENT)
  Comment: "R.POLY.G4: Min Space POLY to Gate for Gate Width > 3.5 must be > 0.45"
  Function: external2
  Inputs: gate_gt_3_5.edgelayer.0001, POLY.polygonlayer.0002
WARNING - 6 spacing violations found.
    Check time = 0:00:00 User=0.00 Sys=0.00 Mem=0.016 GB
\end{lstlisting}

在网络系统上为获得最佳性能，应尽量将所需文件放在本地机器上以加快完成。运行作业的机器本地磁盘应足以容纳 Milkyway 库（或 GDSII / OASIS 文件）、正在创建的 group 文件，以及所有 ASCII 或物理错误文件。

\subsection{运行时间与内存使用报告}
\subsubsection*{Runtime and Memory Use Report}

本节说明摘要文件中针对 IC Validator 运行各阶段报告的运行时间与内存使用。注意：部分处理不在累计内存报告范围内。

\begin{noteBox}
一般而言，峰值或累计内存使用报告仅包含由 IC Validator 工具分配的内存。例如，Milkyway 内存由 Milkyway 库分配，因此不计入 IC Validator 峰值内存或机器累计内存报告。
\end{noteBox}

表~\ref{tab:icv-sum-terms} 给出摘要文件中时间与内存术语的定义。

\begin{table}[htbp]
\centering
\caption{摘要文件中使用的术语}
\label{tab:icv-sum-terms}
\begin{tabular}{@{}lp{0.55\textwidth}@{}}
\toprule
\textbf{术语} & \textbf{定义} \\
\midrule
Time & 墙上时间（wall time） \\
User & 计算所用 CPU 时间 \\
Sys  & I/O 所用 CPU 时间 \\
Mem  & 最大内存使用 \\
\bottomrule
\end{tabular}
\end{table}

以下示例取自摘要报告的不同阶段：

\begin{enumerate}
  \item \textbf{Runset 编译阶段}
  \begin{itemize}
    \item 未使用 runset 缓存时，报告类似：
\begin{lstlisting}
Runset Compile Time=0:01:03 User=61.08 Sys=0.81 Mem=1.065 GB
\end{lstlisting}
    \item 使用 \cmd{-rscache} 启用 runset 缓存时，报告类似：
\begin{lstlisting}
Cached Runset Compile Time=0:00:04 User=0.81 Sys=0.25 Mem=0.203 GB
\end{lstlisting}
  \end{itemize}

  \item \textbf{分布式处理设置阶段}
\begin{lstlisting}
Virtual Machine Init Time=0:00:02 User=1.09 Sys=0.12 Mem=1.065 GB
System Startup Time=0:00:02 User=0.00 Sys=0.52 Mem=1.065 GB
\end{lstlisting}

  \item \textbf{预处理阶段}
  \begin{itemize}
    \item 层次读取与优化：峰值内存报告不计 Milkyway 与 OASIS 内存。
\begin{lstlisting}
Oasis memory used: 27.0M
Milkyway memory usage: 4714.789M
Reading hierarchy Time=0:00:10 User=8.29 Sys=2.87 Mem=0.207 GB
Updating hierarchy Time=0:00:00 User=0.07 Sys=0.00 Mem=0.056 GB
Exploding Time=0:00:00 User=0.44 Sys=0.02 Mem=0.057 GB
\end{lstlisting}
    \item 虚拟 cell（virtual cell）配对阶段：
\begin{lstlisting}
Pairing Iteration 1 Time=0:00:00 User=0.53 Sys=0.04 Mem=0.098 GB
Pairing Iteration 2 Time=0:00:00 User=0.54 Sys=0.05 Mem=0.097 GB
Pairing Iteration 3 Time=0:00:00 User=0.51 Sys=0.04 Mem=0.096 GB
Pairing Iteration 4 Time=0:00:00 User=0.45 Sys=0.03 Mem=0.098 GB
Pairing Time=0:00:01 User=2.03 Sys=0.16 Mem=0.098 GB
Combined VCELL Time=0:00:02 User=2.67 Sys=0.23 Mem=0.098 GB
Post VCell Time=0:00:00 User=0.12 Sys=0.00 Mem=0.060 GB
Prototype Pass 2 Time=0:00:00 User=0.04 Sys=0.00 Mem=0.061 GB
\end{lstlisting}
    \item Regioning 阶段：
\begin{lstlisting}
Determine Region Time=0:00:00 User=0.02 Sys=0.00 Mem=0.059 GB
\end{lstlisting}
    \item Create Layer 阶段：
\begin{lstlisting}
Create Layer Setup Time=0:00:00 User=0.06 Sys=0.01 Mem=0.058 GB
Hierarchy Cleanup Time=0:00:01 User=0.89 Sys=0.04 Mem=0.059 GB
\end{lstlisting}
  \end{itemize}

  \item \textbf{命令执行}
  \begin{itemize}
    \item 命令统计报告：
\begin{lstlisting}
Check Time=0:00:00 User=0.00 Sys=0.00 Mem=0.017 GB
\end{lstlisting}
    \item 读取错误层次：
\begin{lstlisting}
Elapsed Time=0:00:00 User=0.08 Sys=0.00 Mem=0.024 GB
\end{lstlisting}
    \item \cmd{text\_net()}：
\begin{lstlisting}
Total Text Time=0:00:00 User=0.05 Sys=0.00 Mem=0.064 GB
\end{lstlisting}
    \item \cmd{extract\_devices()}：
\begin{lstlisting}
Device Connect Time=0:00:04 User=1.41 Sys=2.84 Mem=0.147 GB
\end{lstlisting}
    \item \cmd{netlist()}：
\begin{lstlisting}
Netlisting Time=0:00:03 User=1.63 Sys=0.68 Mem=0.069 GB
\end{lstlisting}
    \item \cmd{write\_milkyway()}：峰值内存报告不计 Milkyway 内存。
\begin{lstlisting}
Milkyway memory usage: 4129.207M
Compression Time=0:05:07 User=275.34 Sys=27.80 Mem=3.529 GB
Writing Time=0:10:32 User=516.55 Sys=47.36 Mem=3.812 GB
\end{lstlisting}
    \item \cmd{write\_oasis()}：峰值内存报告不计 OASIS 内存。
\begin{lstlisting}
Merging input layout from <input OASIS>
Merge Time=0:00:00 User=0.00 Sys=0.00 Mem=0.025 GB
Writing Time=0:00:00 User=70.00 Sys=0.00 Mem=0.025 GB
\end{lstlisting}
    \item \cmd{write\_gds()}：
\begin{lstlisting}
Writing Time=0:00:00 User=0.00 Sys=0.00 Mem=0.025 GB
\end{lstlisting}
  \end{itemize}

  \item \textbf{主机摘要}
  \begin{itemize}
    \item 每台主机的总运行时间与峰值内存（不计读写 Milkyway 与 OASIS 的内存）：
\begin{lstlisting}
Host number 0 total check Time=0:04:40 User=182.60 Sys=37.42 Mem=0.791 GB
\end{lstlisting}
    \item 每台主机含预处理的总运行时间与峰值内存（不计读写 Milkyway 与 OASIS 的内存）：
\begin{lstlisting}
Host number 0 total check and preprocess Time=0:04:56 User=195.38 Sys=40.62 0.791 GB
\end{lstlisting}
  \end{itemize}

  \item \textbf{错误写入摘要}
\begin{lstlisting}
Overall error storage time: User=0.06 Sys=0.11 Mem=0.014 GB
Generation Time=0:00:00 User=0.01 Sys=0.00 Mem=0.0147 GB
\end{lstlisting}

  \item \textbf{总体内存与运行时间报告}
  \begin{itemize}
    \item 每台机器的累计峰值内存（不计读写 Milkyway 与 OASIS）：
\begin{lstlisting}
IC Validator Peak Machine Memory Report
igcae209     : Peak = 2.061 GB
igcaew008    : Peak = 3.813 GB
igcae191     : Peak = 1.272 GB
igcae177     : Peak = 1.396 GB
igcae188     : Peak = 1.403 GB
igcae205     : Peak = 1.269 GB
igcae178     : Peak = 1.399 GB
igcae174     : Peak = 1.404 GB
\end{lstlisting}
    \item 所有阶段与函数的总体运行时间与峰值内存（不计读写 Milkyway 与 OASIS）：
\begin{lstlisting}
Overall engine Time=0:06:12 Highest command Mem=0.791 GB
\end{lstlisting}
    \item Master 内存使用（不计入峰值内存报告）：
\begin{lstlisting}
Overall Master Mem=0.246 GB
\end{lstlisting}
  \end{itemize}
\end{enumerate}

\subsection{操作系统}
\subsubsection*{Operating Systems}

IC Validator 支持多种硬件架构与操作系统。受支持操作系统列表见安装指南：\\
\url{http://www.synopsys.com/support/licensing-installation-computeplatforms/compute-platforms/release-specific-support.html}

\subsection{系统限制}
\subsubsection*{System Limits}

若系统限制设置不当，IC Validator 运行中可能出错。例如，下列命令为 Linux 上的 IC Validator 运行设置合适的系统限制（数值单位为 kilobytes）：
\begin{lstlisting}[style=shell]
unlimit filesize
unlimit datasize
limit descriptors 1024
limit stacksize 8192
\end{lstlisting}

\begin{noteBox}
\cmd{descriptors} 设置控制运行期间可同时打开的文件数量。
\end{noteBox}

要验证主机上的限制，请运行 shell 的 \cmd{limit} 命令。

\subsection{网络环境}
\subsubsection*{Networked Environments}

在网络环境中访问文件（例如使用远程 group 目录）可能降低性能。您可能观察到 IC Validator 报告的墙上时间（wall time，总体时间）与总 CPU 时间（user time + system time）之间存在很大差异。

\begin{noteBox}
这些差异也可能由下列原因造成：
\begin{itemize}
  \item 同一系统上运行的其他进程；
  \item 系统开销（system overhead）。
\end{itemize}
\end{noteBox}

NFS 挂载的读/写传输大小（\cmd{rsize} 与 \cmd{wsize}）应至少为 32768。例如，对静态挂载设置：
\begin{lstlisting}[style=shell]
% mount -t nfs -o rsize=32768,wsize=32768 server:/path /mnt/nfs
\end{lstlisting}
